\section{CONCLUSIONES}

Con el simulador fue posible realizar los seis experimentos. Tres de ellos son comportamientos ampliamente conocidos, en los que es posible saber exactamente cuál será el siguiente estado: (1) el juego de la vida, (2) el \textit{sandpile} y (3) la hormiga de Langton. Con ellos se pudieron definir configuraciones de los autómatas celulares neuronales y, más específicamente, de la red neuronal, que facilitan o permiten inferir un comportamiento.


Con el juego de la vida se demostró la capacidad inicial del sistema para simular sistemas complejos. El juego de la vida requiere una de las redes neuronales más pequeñas por la poca información que necesita: 9 neuronas en la capa de entrada, 1 capa oculta de 4 neuronas con ReLU y 2 neuronas con función Linear en la capa de salida. Además, este experimento muestra la relación entre las reglas arbitrarias y las semitotalísticas, pues un entrenamiento incompleto, variando la pérdida objetivo, permite que la red aprenda otra regla.

Se observa una estabilidad en la regla B3/S23, pues al variar la pérdida objetivo, únicamente se aprenden 2 comportamientos: la regla 0, en la que cualquier patrón se vuelve 0, o la misma regla. Lo contrario ocurre con otras reglas, como B3/S238, B38/S238 o B3/S236, que generan vida adicional y que, al variar su pérdida objetivo, sí llegan a aprender otras reglas. Esto puede deberse a que la red neuronal se inicializa generando ceros y, como casi nunca recibe los patrones de baja frecuencia, sus salidas para ellos se mantienen en ese valor. También cabe la posibilidad de que esta regla sea más estable que las demás.

Finalmente, este primer experimento muestra la importancia de la vecindad al entrenar: con una vecindad más pequeña no se puede aprender el comportamiento, pero una vecindad más grande no impide replicarlo.

Con la pila de arena de Bak–Tang–Wiesenfeld se tienen 17 estados. Este experimento tuvo dos aportaciones principales. La primera fue reducir la tasa de aprendizaje con \textit{plateau} para evitar oscilaciones en la pérdida y llegar a la pérdida objetivo de manera uniforme. La segunda fue usar codificación \textit{one-hot} para que los estados no tengan una relación lineal entre sí. Sin embargo, la red no pudo extrapolar el comportamiento de una célula generadora a varias, aunque sí generalizó parcialmente hacia generadoras contiguas.

El último de estos tres experimentos sirvió para demostrar la capacidad de manejar múltiples estados y escasez de información. La hormiga de Langton únicamente cambian 2 células en cada paso y el resto se mantiene constante. A pesar de esta escasez de información, sin requerir preprocesamiento adicional, el autómata celular neuronal demostró su capacidad de extraer información de comportamientos poco representados, con una alta efectividad incluso usando 10 colores (50 estados). Con una sola hormiga evolucionó sin errores hasta el paso 12,455, donde la hormiga se divide en dos al aplicar una regla de colisión aprendida de los datos con dos hormigas.

Queda demostrado que el autómata celular neuronal puede replicar algunos comportamientos complejos, especialmente de configuraciones bien conocidas, pero con restricciones. La mayor limitación es la representación de escenarios.

Esta demostración dio pie a los siguientes tres experimentos, que pusieron a prueba la capacidad de replicación en entornos reales: (1) la propagación de fuego con vista cenital (FLAME~3), (2) el movimiento de tormentas (SEVIR) y (3) el crecimiento de moho en una caja de Petri.

La propagación del fuego se puede replicar, pero existen múltiples elementos que no se consideran y que limitan la capacidad de replicación, que al ignorarlas todas, la abstracción dentro del simulador pierde tanto contexto del el sistema real. El autómata celular neuronal reproduce el avance del fuego, pero la velocidad y la dirección dependen exclusivamente del conjunto de datos con el que fue entrenado. Además, este sistema se vio limitado por tener únicamente dos grabaciones pequeñas, los comportamientos aprendidos dependen mucho de las condiciones en las que se grabaron: o el sistema se calienta de manera uniforme y tiene un salto conforme se acerca la barrera de fuego o el sistema está frío y las temperaturas solo aumentan cuando llega la barrera de fuego. 

Por otro lado, la simulación del movimiento de tormentas presenta otro inconveniente. El sistema es capaz de replicar el movimiento, pero únicamente en las secuencias que se mueven en una sola dirección. Como se muestra en el último experimento, una alternativa es agregar el ruido al estado de cada célula sin cambiar dicha definición. Por último, lo que hace que este sistema no sea apto para simularse con un autómata celular neuronal es que no es un sistema cerrado: existe factores externos que intervienen en su comportamiento. Por esta razón se propone el último experimento, que consiste en grabaciones del moho \textit{Physarum polycephalum} en una caja de Petri, un sistema cerrado que tiene una configuración inicial.

El autómata celular neuronal sí aprende el crecimiento del moho a partir de su evolución: el plasmodio crece desde el inóculo, la avena lo guía y no aparece en zonas vacías. El paso de tiempo resultó ser lo más importante: la coincidencia del plasmodio sube de 0.120 con pasos de 30 minutos a 0.223 con pasos de 6 horas. Elegir siempre el estado más probable elimina la parte aleatoria del crecimiento, por lo que la simulación crece en manchas sólidas y sin movimiento. Agregar un valor de ruido al estado de cada célula permite reproducir esa parte aleatoria con una regla determinista que el simulador ejecuta sin cambios: en un paso, la red acierta el estado visible del 99.17\% de las células y el ruido del 99.98\%, y con ella aparecen las ramas y el movimiento interno. Lo que la red no logra es formar las venas largas que unen la avena ni saber cuándo dejar de crecer, ya que probablemente dependen de lo que ocurrió durante horas, mientras que la red solo conoce el estado actual.

La hipótesis se cumple en los sistemas cerrados. El autómata celular neuronal infiere una regla local que reproduce la evolución con la que fue entrenado. No se cumple en sistemas abiertos como las tormentas, cuya evolución depende de factores externos que no están en el espacio, y se cumple parcialmente en la propagación del fuego, limitada por las variables que no se registran y por la cantidad de datos. En todos los casos, reproducir la evolución de entrenamiento no garantiza haber extraído la función de transición completa: ante patrones que no aparecen en los datos de entrenamiento surgen errores.